%% GENERATED by scripts/paper/build_prose.py from scripts/paper/prose.yaml (the prose) and the analysis
%% register (the options blocks) -- edit the prose there, or here once you stop regenerating.

% ======== q-trust-unchecked ========

\subsection{You trust what you cannot check}\label{sec:qual:trust-unchecked}

The most articulate operators described trusting the Tactical Assistant \emph{because} they could not assess its plan, while a strategy card was simple enough to judge at a glance, and so was judged, often harshly. Under time pressure checking shrank to a glance, and the minority who did check earned their trust by comparing the assistant's plan with their own.

% --- q-trust-unchecked#cant-check (pass 5: for the paper) ---

\phantomsection\label{sec:qual:trust-unchecked:cant-check}
A strategy card shows a few numbers per drone type beside the mission's requirements, so its oddities are visible; a tactical plan is a web of routes on a map. Participants drew the consequence themselves: they took the tactical plan on trust because they could not see how far it was from optimal, and suspected that was part of why it seemed so good. This is their own account of two of the quantitative results, the trust gap between the tiers (Section~\ref{sec:res:tier-trust}) and the visible and invisible flaws (Section~\ref{sec:res:flaw-odd}), and it suggests the gap is partly one of checkability rather than quality.

\begin{quote}\small
``The tactical assistant, I never really saw it do anything wrong, but that's basically because I think the problem is so complex [that] seeing a more optimal solution is really hard. So I wonder if it's one of those things where these are based on the fact that I actually can't perceive the difference between its performance and the optimal.''\par\hfill--- P22 (talk during the questionnaire survey1 0:14:45)
\end{quote}

\begin{quote}\small
``I will be honest with you. I don't think I was paying as much attention to the suggestions of the tactical assistant as I could have been. Just purely because the tasks were often very multi-drone. I think I would have committed a lot of extra thought-space to assigning all of the drones individually to those. So I definitely inadvertently trusted the tactical [agent]. I think it's worth saying any assessment of the tactical system would probably be disingenuous just because of my lack of engagement with it. But the strategic assistant, I think on the most part, I have absolutely no problems with. There was that one issue, but that's okay because I noticed it because the display is straightforward and it's clear enough that you can as a human spot discrepancies if they arise.''\par\hfill--- P19 (talk during the questionnaire survey1 0:17:06)
\end{quote}

\begin{quote}\small
``i felt like i took a bigger pinch of salt for the strategic one because i felt like i could assess the quality of it more immediately whereas with the tactical assistant, I couldn't accurately and rapidly assess how well it was performing. So I just kind of said, OK, I have to take your word for it, essentially.''\par\hfill--- P22 (talk during the questionnaire survey2 0:26:53)
\end{quote}

% --- q-trust-unchecked#glance (pass 5: extra: may make it in) ---

\phantomsection\label{sec:qual:trust-unchecked:glance}
Nobody described careful checking under pressure. They described a quick look, and accepted the trade because penalty accrued while they checked; one participant set out the expected-value argument directly.

\begin{quote}\small
{\footnotesize\textit{Researcher:} Am I correct in saying you should be using the agents a little bit more often now?}\par\smallskip
``Whatever lack of optimality that comes from using the AI agents is more than compensated for in their speed. Or at least in their speed of decision making. I also feel like the longer that this session has gone on, the less oversight that I've been giving it, and the more I've been focusing on the strategic layer and the incoming missions.''\par\hfill--- P31 (think-aloud s2 0:20:58)
\end{quote}

\begin{quote}\small
{\footnotesize\textit{Researcher:} So you say you especially trust it with a large mission. You mean you're happy to rely on it because of the speed?}\par\smallskip
``Yeah. I think even if the tactical agent gave me a dodgy suggestion and the mission failed and I lost points, I probably would have gained more points from the time saved''\par\hfill--- P11 (debrief 0:29:11; includes the researcher's words)
\end{quote}

% --- q-trust-unchecked#earned (pass 5: for the paper) ---

\phantomsection\label{sec:qual:trust-unchecked:earned}
Where trust was earned, it was by comparison with what the operator would have done themselves.

\begin{quote}\small
{\footnotesize\textit{Researcher:} So I would start with just how did you find that? Anything immediately that you would comment on?}\par\smallskip
``I think I developed a stronger opinion on how much to trust each of them. By, I guess, comparing them to what I would have done. And I mean, I understood the task more as I went. So the tactical one seemed to be doing what I was doing. So quite reliably. And obviously the strategic one wasn't. So I guess it was very much comparing them to what I'd do. And then letting the tactical one get on with it.''\par\hfill--- P32 (debrief 0:23:45)
\end{quote}

\begin{quote}\small
{\footnotesize\textit{Researcher:} The degree to which you trusted the decisions the AI was making, uh, and how that progressed throughout the study?}\par\smallskip
``I mostly trusted it, although, uh, I did kind of want to test it more at the beginning, like comparing it, like doing a few manual ones myself to see, like, whether it would be better, faster. I guess it was kind of hard to keep track of which tasks were being finished when, but I noticed that the ones that failed were normally cause I'd just done a manual one. So I guess that made me want to use the AI more often.''\par\hfill--- P29 (debrief 0:25:59)
\end{quote}

\begin{options}
\optshead{Options for this finding}

\optsec{Trusted because it cannot be checked} {\footnotesize [pass 5: for the paper]}

\noindent Your note: A nice few quotes that we can use to make a point.

\optsec{Checking shrinks to a glance} {\footnotesize [pass 5: extra: may make it in]}

\noindent Your note: Less significant but might be worth a mention

\noindent Highlighted quotes not used above:\begin{quotelist}
\item ``Speed is the name of the game. So sometimes in these ones, I guess making a decision earlier is beneficial because at some point the other ones are going to pop up.''\,---\,P22 (think-aloud s1 0:04:20, includes the researcher's words)
\end{quotelist}

\optsec{Trust earned by comparison, and doubts that came with practice} {\footnotesize [pass 5: for the paper]}

\noindent Your note: Minimal but nice
\end{options}

\clearpage

% q-process (Trust in the process, not the answer: what operators thought the assistants were) left out: your verdict this pass was prune.

% ======== q-who-decides ========

\subsection{Who decides, and who answers for it}\label{sec:qual:who-decides}

Almost everyone wanted a human to keep the decision, and the reasons were moral as much as practical: accountability for a life not saved, a human connected to the operation, and the plan held in someone's head. People were candid that the human they wanted in the loop might not be able to check much.

% --- q-who-decides#accountability (pass 5: for the paper) ---

\phantomsection\label{sec:qual:who-decides:accountability}
The clearest statement of the moral reason is about accountability: even if an AI saves more lives on average, when it fails to save one there is no one to answer for it. Against it sits the plain value of speed in an emergency.

\begin{quote}\small
{\footnotesize\textit{Researcher:} But one of my questions is if in five years' time this was much more sophisticated, far more complicated, how would you feel about search and rescue scenarios being conducted with AI systems of this type in this kind of setup?}\par\smallskip
``Yeah, I feel like that's a very interesting question for more of a moral reason, because if you look at it entirely from the perspective of your option A, option B, and option A saves more lives, that always seems like a good thing, but taking it ... in a slightly less simplified way, you've got to make sort of moral judgments about how these resources are being allocated, and if someone's life was like failed to be saved, under a completely human system, you'd have some sort of like, not responsibility, what's the word I'm looking for? [interviewer suggests the word 'accountability']. Accountability, thank you! Of who was on the job that failed to save your life, whereas if it's an AI system, even if it performs on average better than a human, that accountability just isn't there.''\par\hfill--- P31 (debrief 0:27:53; includes the researcher's words)
\end{quote}

\begin{quote}\small
{\footnotesize\textit{Researcher:} How do you feel about people using systems like this in, um, high risk scenarios?}\par\smallskip
``I think definitely with emergency scenarios, it's best to just get something going, even if it's not the most perfect plan. And AI lets you do that quite good, like quite quickly.''\par\hfill--- P29 (debrief 0:27:05; includes the researcher's words)
\end{quote}

% --- q-who-decides#loop (pass 5: for the paper) ---

\phantomsection\label{sec:qual:who-decides:loop}
Others wanted a human in the loop whatever that human could check, and one participant framed it organisationally: measure where people and assistants each add value, and assign the tiers accordingly.

\begin{quote}\small
{\footnotesize\textit{Researcher:} If we took this five years in the future, and systems of this type were being used in real search and rescue scenarios with more complicated AIs and operators being assisted, how do you feel about that? And do you see any potential for good and for bad of doing that sort of thing?}\par\smallskip
``I think if I were being organisationally strategic about that, it seems like it's really helpful. Assuming that you're human constrained (if you've got loads of spare humans, I'd want to see if the AI was actually better) but assuming that you've got less humans than you ideally would have, it feels like a good accelerator. I'd want the analytics on, I guess, like drone utilisation and mission success. I feel like that would be a good place to be making these decisions. And should I be, for example, in this case where I think the strategic one was less good than the tactical, if I were being really strategic, I'd be saying: 'okay, people are going to do the strategic stuff, and we're just going to trust it for the tactical'. So I guess very much doing that optimisation around: 'is it the agents or the humans who are adding the most value, and where do they add the most value', and then have them allocate them... I'd want to understand, well, just because I like to understand everything, I'd want to understand empirically actually how it was performing, who should be doing what.''\par\hfill--- P32 (debrief 0:26:10; includes the researcher's words)
\end{quote}

% --- q-who-decides#ownership (pass 5: for the paper) ---

\phantomsection\label{sec:qual:who-decides:ownership}
For some the reason to keep the work was not safety but connection to the task, and pride in their own solution.

\begin{quote}\small
``I don't want to be reduced to pressing allocate, suggest, deploy, because otherwise I might as well be in the next room. And I think the small missions give me the chance to do that more. And I'm like, 'oh, I can just do that. Don't worry. I'll do that'. Done. And I feel satisfaction.''\par\hfill--- P11 (debrief 0:30:55)
\end{quote}

\begin{quote}\small
``it would be a dreadful task without any of the tools, especially on the tactical side. It's just a lot of laborious dragging, holding shift at the right time.... And I guess I do feel I want to use (empahsis)my(/emphasis) solution, and I don't want to like, oh, you know, of course the computer thinks it's better. Or like personify the computer and feel some sort of grudge to hold it.''\par\hfill--- P11 (debrief 0:26:55; includes the researcher's words)
\end{quote}

\begin{options}
\optshead{Options for this finding}

\optsec{Ownership, and the opposite view} {\footnotesize [pass 5: for the paper]}

\noindent Your note: This is one of the better ones
\end{options}

\clearpage

% ======== q-failure ========

\subsection{Failure without feedback}\label{sec:qual:failure}

Drone failures were random and uniform, and nothing on screen said so or said how likely they were. Operators felt the lack, built their own theories about where failures happened, and blamed themselves.

% --- q-failure#blind (pass 5: for the paper) ---

\phantomsection\label{sec:qual:failure:blind}
Several wanted failure rates built into the allocation, so that the reserve could be sized to the risk rather than padded.

\begin{quote}\small
``I think in a in a real world scenario where if I imagine I'm a drone swarm operator you will have some numbers presumably about 'so what's the likelihood for drone failure' right? Based on historic data, and that will presumably bias (or inform, not bias) but inform your selection if I'm like okay the lifter drone is much heavier therefore there's more wear and tear on it or whatever it's more likely to fail down that means that my redundancy strategy might change from 'send one additional drone for the lifter' towards 'send 20 percent more more lifters than you need''\par\hfill--- P10 (debrief 0:28:53; includes the researcher's words)
\end{quote}

\clearpage

% ======== q-hierarchy ========

\subsection{The hierarchy as one plan}\label{sec:qual:hierarchy}

Operators did not experience two independent assistants. They described one plan that starts at the strategic step: allocate too little at the top and the shortage cascades down; allocate generously and there is only one sensible way to assign the drones; build the allocation yourself and the tactical plan is already in your head.

% --- q-hierarchy#cascade (pass 5: for the paper) ---

\phantomsection\label{sec:qual:hierarchy:cascade}
Several placed the responsibility for failures with the strategic tier, since only drones sent at the start can absorb them; others weighed the same step from the other side, holding drones back for missions not yet arrived.

\begin{quote}\small
``Basically, it did everything. If failure rate of drones is unrelated to my assignment, so rate of failure as opposed to how many drones I send out, then all things considered, the tactical assistant is probably the best. The tactical assistant completed all the jobs and we didn't end up in any kind of a situation where we didn't have enough drones if it hadn't been for random drone failures. And the only way to remedy drone failures is to have sent enough drones at the beginning, which I think is the strategic assistant's prerogative.''\par\hfill--- P19 (talk during the questionnaire survey1 0:18:53)
\end{quote}

\begin{quote}\small
{\footnotesize\textit{Researcher:} You found that one more demanding?}\par\smallskip
``I'm trying to be more defensive in the strategy, so I'm trying to reserve more and more drones in the base, rather than send them out to finish the task quickly, because I cannot predict when the next task will appear. So luckily, there's only one task is holding for a long time, only one there, and another problem is [that] such a choice will make the drones in the case becomes not sufficient for that case. So if I only send out one drone for one specific color, and it fails, then I have to abort that mission, and I only abort one of them, because that happened, and the rest is still able to finish, so I think this, such a deployment is still all right.''\par\hfill--- P28 (talk during the questionnaire survey2 0:28:05)
\end{quote}

% --- q-hierarchy#one-plan (pass 5: for the paper) ---

\phantomsection\label{sec:qual:hierarchy:one-plan}
The link between the tiers in the logs (Section~\ref{sec:res:size-decides:tactical}) was described from both sides: a generous strategy card makes a tactical suggestion the obvious next step, while a hand-built allocation already lays out a plan the operator does not want disrupted. Whether one plan held across both tiers is the right ideal is a question the practitioner's account takes up, and answers differently (Section~\ref{sec:qual:practitioner:pt-silos}).

\begin{quote}\small
``it's better to use the autonomy on the tactical level on top of the strategic''\par\hfill--- P10 (debrief 0:25:46; includes the researcher's words)
\end{quote}

\begin{quote}\small
``[Having manually set a strategic allocation] I prefer to [tactically] set the drones myself. Because I'm not sure how exactly it will be allocated [if I use the agent]. Because I have something in mind when I consider the allocation... When I make the selection for the drones, I kind of lay out a plan in my mind initially. So I feel like it's going to disrupt my thought process if I try to go around that.''\par\hfill--- P21 (think-aloud s1 0:14:16; includes the researcher's words)
\end{quote}

\begin{quote}\small
``Actually, they might work better if I use the AI agent for the selection and that as well, because. Yeah, because when I assume when it makes a selection, when it tries to make the drone allocation, it follows certain patterns, which could then be adopted by the tactical agent afterwards. And if I try to think of something else, so I'm going to split them like that and change them in a certain way. I'm not sure if the tactical agent will be able to match my thought process.''\par\hfill--- P21 (talk during the questionnaire between1 0:20:36)
\end{quote}

% --- q-hierarchy#cross-mission ---

\phantomsection\label{sec:qual:hierarchy:cross-mission}
A strategy card considers one mission. Operators overrode the cards to act on what the cards could not see: keeping drones back for other missions, and weighing which missions could afford to wait.

\begin{quote}\small
{\footnotesize\textit{Researcher:} Could you just summarize again, just in a sentence or two, how you were approaching each one, what scenarios you felt you were using each one?}\par\smallskip
``I felt like sometimes I could find like a middle ground that had like maybe a little bit of gambling in there, essentially, that you wouldn't have the necessary coverage. But then also then that meant that you'd keep some more in reserve. So it wasn't just fully because it's cautious, it was also because it was actually to keep some in reserve for other tasks, because, yes, you're covering this task specifically very well. But that doesn't mean that then you actually have enough coverage for the, for the rest of your tasks. So I felt that often I could find, I could strike a middle ground. So I felt comfortable assessing and then overriding the strategic assistant.''\par\hfill--- P22 (debrief 0:28:28)
\end{quote}

\begin{quote}\small
{\footnotesize\textit{Researcher:} Now that I've had a little bit more time, I'm just thinking if my opinions have changed on this slightly. And for what reason?}\par\smallskip
``If you wanted to send as few drones as possible, you could send four fast drones, one lifter, one camera. Whereas even the conservative recommendation recommended sending two of the lifters, and the cameras in case of drone failure. When I was in that situation, I had a lot of other pending tasks, so I went more conservative than the conservative recommendation.''\par\hfill--- P31 (talk during the questionnaire survey2 0:23:49; includes the researcher's words)
\end{quote}

\begin{options}
\optshead{Options for this finding}

\optsec{The cross-mission view the cards lack}

\noindent Highlighted quotes not used above:\begin{quotelist}
\item ``And the problem is we don't know whether a extremely critical job is going to show up later.''\,---\,P19 (think-aloud s1 0:09:07)
\end{quotelist}
\end{options}

\clearpage

% q-real-ops (What a real operation would need) left out: your verdict this pass was prune.

% ======== q-new-jobs ========

\subsection{An assistant for what to open next}\label{sec:qual:new-jobs}

Asked what the assistants should do, participants rarely asked for more automation of the same kind. Two asked for something neither assistant does: help deciding which waiting mission to open next.

% --- q-new-jobs#queue-order (pass 5: for the paper) ---

\phantomsection\label{sec:qual:new-jobs:queue-order}
Neither assistant helps with the first decision of every allocation, which waiting mission to open. With two or more missions waiting (411 decisions), operators opened the oldest well above chance (49\%, against 34\% by chance) and leaned towards the more critical (odds $\times 1.45$ per category above the waiting missions' average, $p = .02$), but not towards the missions they could actually cover: whether the reserve could supply a mission's drones made no difference ($\times 1.04$, $p = .79$), and at 41\% of these decisions the oldest waiting mission could not be covered. An ordering of the kind one participant proposed (missions that can be done first, then by criticality) would have picked a different mission at 59\% of decisions.

\begin{options}
\optshead{Options for this finding}

\optsec{An assistant for what to open next} {\footnotesize [pass 5: for the paper]}

\noindent Key numbers:\begin{itemize}[leftmargin=1.2em,itemsep=0pt]\footnotesize
\item Decisions with 2+ missions waiting: 411 (33 people). Opened the oldest 49\% (chance 34\%); when one waiting mission was uniquely the most critical, opened it 49\% (chance 33\%)
\item Choice among the waiting missions (conditional logit, person bootstrap, 400 draws): oldest $\times$2.75 [1.97, 4.12], p < .001; newest $\times$1.64 [1.17, 2.32], p < .001; criticality, per category above the set's mean $\times$1.45 [1.09, 2.15], p = .02; coverable from the reserve $\times$1.04 [0.72, 1.58], p = .79; size per 3 drones $\times$0.95 [0.78, 1.11], p = .61
\item The oldest waiting mission could not be covered from the reserve at 41\% of these decisions; some waiting mission could not be covered at 57\%. When some could not, the one opened was coverable 32\% of the time (chance 39\%)
\item P27's order (coverable, then most critical, then oldest) picks a different mission from oldest-first at 59\% of decisions; there operators opened P27's pick 35\% of the time (chance 30\%)
\end{itemize}
\end{options}

\clearpage
